\documentclass{article}
\usepackage{graphicx} % Required for inserting images
\usepackage{algorithm} %
\usepackage{algpseudocode}
\usepackage{amsfonts}

\usepackage{tcolorbox}
\newtcolorbox{greenbox}[1]{colback=green!5!white,colframe=green!75!black,fonttitle=\bfseries,title=#1}


\usepackage{amsthm}
\newtheorem{definition}{Definition}
\newtheorem{theorem}{Theorem}
\newtheorem{lemma}{Lemma}

\newtheorem*{remark}{Remark}

\usepackage{amsmath}
\DeclareMathOperator*{\argmin}{argmin} % no space, limits underneath in displays



\setlength{\parindent}{0pt} % for getting rid of first line indentation


% ================================================
% =====================  BEGIN DOCUMENT 
% ================================================

\title{Forecasting}
\author{Juan Pablo Bello Gonzalez}
\date{January 2024}

\begin{document}



\maketitle

\section{Statistics stuff}
\begin{itemize}
\item $Cov(X_i, X_j) = 0$ for uncorrelated RVs
\item for two RVs, Independence is a stronger condition than non-correlation


\end{itemize}

\section{Introduction to Forecasting}

\begin{definition}
Time series: A sequence of random variables indexed by time $\{ X_t : t\in T \subset \mathbb{N}\}$. It is a special case of a stochastic process. 
\end{definition}

\begin{definition}
Stochastic process: a collection of random variables (they don't need to have the same distribution!) indexed by a set $I$. $I$ could be continuous or discrete.
\end{definition}

We often want to model time series data with \textit{non-stationary trends}:
\begin{itemize}
\item Seasonality
\item trending mean
\item Non-constant variance (fanning in our out pattern)
\end{itemize}

Ideally given a time series $\{X_1, X_2, ... X_n\}$ we would know the \textit{joint distribution} of such vectors. 
\[ F(x_1, x_2, \dots x_n) = P(X_1, X_2, \dots X_n)\]

This is because individual (marginal) distributions of $X_i$ don't tell you about the combined behaviour.

We want to study a stochastic process as a mathematical which is a sequence of RV's with some predictable distribution ie not just a vector of n random and unrelated numbers. As such, the existence of a stochastic process should always be questioned. Kolmogorov's theorem tells us about the existence of  discrete time stochastic process build from $X_i's$ as long as the $X_i$ satisfy the two xonsistency conditions. 

If all finite dimensional distributions exist, and they are consistent with each other under marginalization, then the stochastic process exists.  Then \textbf{we can use finite dimensional models ot define a general process}.*****

All finite dimensional distributions of $X_1, X_2, \dots X_n$ are those of all non-empty subsets of the sequence ie $\{ X_1, X_2, \dots X_n \}^*$ so it is impossible to have all of these in practice since, if we are sampling, we don't even have the distribution of a single $X_i$ and must do with estimates. 

Fortunately, in most application we can get by with knowing $E[X_t], E[X_t^2], E[X_tX_j]$. Which is equivalent to knowing $E[X_t], Var(X_t), Cov(X_t, X_j)$ for \textbf{ALL} indices.

\begin{definition}
\textbf{A time series model} for the observed data $\{x_t\}$ is a specification of the joint distributions or more often in practice, the means and covariances of a sequence of random variables $\{X_t\}$ of which $\{x_t\}$ is a realization. 
\end{definition}


\subsection{Zero-mean time series models}
\subsubsection{i.i.d Noise}
\begin{definition}
\textbf{iid Noise} a sequence $\{X_i : i = 1,2\dots n\}$ with: 
\begin{itemize}
\item  independence property is satisfied: joint = product of marginals.
\item  $E[X_i] = 0 \forall i$

\end{itemize} 
\end{definition}

As is the case with a discrete time stochastic process with independent components, the memory-less property is satisfied.  
\begin{proof}
\[ P(X_{n+h} \leq x | X_1 \leq x_1, \dots, X_1 \leq x_n ) = \frac{P(X_{n+h} \leq x)  F(x_1, \dots x_n)}{F(x_1, \dots x_n)}\]
\end{proof}
Note that in the nominator above, the probability of the intersection of events is equal to the individual probabilities multiplied by independence. 

Also note how we can have the individual $X_i$ be continuously distributed but discrete stochastic process (indexed by the naturals)

\subsubsection{White Noise }
Sequence of uncorrelated random variables $\{X_1, X_2, \dots \}$ with zero mean and finite variance. Denoted $X_i \sim WN(0, \sigma)$
\begin{itemize}
\item $E(X_i)=0$ 
\item $Var$
\item $Cov(X_i, X_j) = 0$ for uncorrelated 
\end{itemize}


\subsubsection{Random Walk}
\begin{definition}
A R.V which is a \textbf{series} of \textbf{iid noise OR White noise} distributed R.V.s $S_t = \sum_{i=1}^t X_i$. Usually starting at zero. 
\end{definition}

What is the relationship between \textit{white noise} and \textit{iid noise}?
Independence is a stronger condition than non-correlation. For independence we need to know joint distributions but for non-correlation we don't. So we'd rather work with white noise. 

\section{Classical decomposition}
Fitting models with trend and seasonality. 

In regression we are always modelling $Y|X$ and not just $Y$ since we are actually using the explanatory variable to fit $Y$.

The general framework is to remove the \textbf{trend}, \textbf{seasonality}, and \textbf{non-constant variance} patterns in the data  so that \textit{hopefully} we are left with a stationary process. Then model the detrendized and deseasonalized data. 










\end{document}
